The solution of the Poisson equation (\ref{eq:PoissonEquation}),
\begin{equation}
\nabla^2 \phi (t, \textbf{x})= 4\pi G \rho (t, \textbf{x}),
\end{equation}
can be written formally in terms of a Green function. 

The standard mathematical definition of the Green's function for a linear differential operator $\hat{L}$ is the solution to $\hat{L}\mathcal{G} = \delta$, where $\delta$ represents the Dirac delta function for all relevant coordinates. For the gravitational system, the coordinates are time $t$ and space $\textbf{x}$. Therefore, our defining equation is
\begin{eqnarray}
    \nabla^2 \mathcal{G}(t, \textbf{x}; t', \textbf{x}') = \delta(t - t') \delta^{(3)}(\textbf{x} - \textbf{x}').
\end{eqnarray}

However, it is important to note that  the differential operator $\nabla^2$ contains only spatial derivatives, so it does not act on the time coordinate (remember that in Newtonian mechanics, gravitational influence propagates instantaneously). Because time acts merely as a parameter rather than a dynamic variable in this operator, the Green's function separates completely into a temporal part and a spatial part
\begin{equation}
    \mathcal{G}(t, \textbf{x}; t', \textbf{x}') = \delta(t - t') \mathcal{G}_s(\textbf{x}, \textbf{x}').
\end{equation}
Substituting this back into the definition, yields the spatial Green's function equation
\begin{equation}
    \nabla^2 \mathcal{G}_s(\textbf{x}, \textbf{x}') = \delta^{(3)}(\textbf{x} - \textbf{x}').
\end{equation}
Due to translation invariance of free space, $\mathcal{G}_s$ depends only on the relative distance $\textbf{R} = \textbf{x} - \textbf{x}'$, so the equation can be written as 
\begin{equation}
    \nabla^2 \mathcal{G}_s(\textbf{R}) = \delta^{(3)}(\textbf{R}).
\end{equation}

While this can be solved using Fourier transforms, a highly elegant and rigorous physical method is to use Gauss's Divergence Theorem. Thus, we integrate both sides of the spatial equation over a spherical volume $V$ of radius $R=|\textbf{R}|$ centered at the origin,
\begin{equation}
    \int_V \nabla^2 \mathcal{G}_s(\textbf{R}) \, d^3R = \int_V \delta^{(3)}(\textbf{R}) \, d^3R.
\end{equation}

From the definition of the 3D Dirac delta function, the volume integral over a space containing the origin is exactly $\int_V \delta^{(3)}(\textbf{R}) \, d^3R=1$. Then,
\begin{equation}
    \int_V \boldsymbol{\nabla} \cdot (\boldsymbol{\nabla} \mathcal{G}_s(\textbf{R})) \, d^3R = 1.
\end{equation}

Applying the Divergence Theorem ($\int_V \boldsymbol{\nabla} \cdot \textbf{A} \, dV = \oint_{\partial V} \textbf{A} \cdot \textbf{da}$), the volume integral becomes a surface integral over the sphere's boundary $\partial V$, 
\begin{equation}
    \oint_{\partial V} \boldsymbol{\nabla} \mathcal{G}_s(\textbf{R}) \cdot \textbf{da} = 1.
\end{equation}

Because the source $\delta^{(3)}(\textbf{R})$ is symmetric, the solution $\mathcal{G}_s(\textbf{R})$ must depend only on the radial distance $R$ and therefore, the gradient $\boldsymbol{\nabla} \mathcal{G}_s(\textbf{R})$ points purely in the radial direction $\hat{\textbf{R}}$,
\begin{equation}
    \boldsymbol{\nabla} \mathcal{G}_s(\textbf{R}) = \frac{d\mathcal{G}_s}{dR} \hat{\textbf{R}}.
\end{equation}
The differential area vector on the sphere's surface is $\textbf{da} = \hat{\textbf{R}} R^2 d\Omega$ (where $d\Omega$ is the solid angle). The integral becomes
\begin{align}
    \oint_{\partial V} \left( \frac{d\mathcal{G}_s}{dR} \hat{\textbf{R}} \right) \cdot \left( \hat{\textbf{R}} R^2 d\Omega \right) = &1 \\
    \frac{d\mathcal{G}_s}{dR} R^2 \oint_{\partial V} d\Omega = &1.
\end{align}
Integrating the solid angle gives the total surface area of a sphere, i.e.  $4\pi$, and hence we obtain
\begin{equation}
    \frac{d\mathcal{G}_s}{dR} (4\pi R^2) = 1 \implies \frac{d\mathcal{G}_s}{dR} = \frac{1}{4\pi R^2}.
\end{equation}
Finally, we integrate with respect to $R$,
\begin{equation}
    \mathcal{G}_s(\textbf{R}) = \int \frac{1}{4\pi R^2} dR = -\frac{1}{4\pi R} + C.
\end{equation}
For an unbounded domain, we demand that the potential vanishes at infinity, i.e. $\lim_{R \to \infty} \mathcal{G}_s(\textbf{R}) = 0$, so the constant of integration is $C = 0$. Replacing $R = |\textbf{x} - \textbf{x}'|$, we get
\begin{equation}
    \mathcal{G}_s(\textbf{x}, \textbf{x}') = -\frac{1}{4\pi |\textbf{x} - \textbf{x}'|}
\end{equation}
and recombining the spatial and temporal parts, the explicit spatio-temporal Green's function is finally
\begin{equation}
    \mathcal{G}(t, \textbf{x}; t', \textbf{x}') = -\frac{\delta(t - t')}{4\pi |\textbf{x} - \textbf{x}'|}. \label{eq:AppendixGreenFunction}
\end{equation}


In order to prove this deduction, we will apply it to find the solution $\phi(t, \textbf{x})$ for the original equation (\ref{eq:PoissonEquation}). The general solution via Green's functions is the integral of the Green's function multiplied by the entire right-hand side source term, i.e. 
\begin{equation}
    \phi(t, \textbf{x}) = \int dt' \int d^3x' \, \mathcal{G}(t, \textbf{x}; t', \textbf{x}') \Big[ 4\pi G \rho(t', \textbf{x}') \Big].
\end{equation}
Substitution of the Green function (\ref{eq:AppendixGreenFunction}) gives 
\begin{equation}
    \phi(t, \textbf{x}) = \int dt' \int d^3x' \left( -\frac{\delta(t - t')}{4\pi |\textbf{x} - \textbf{x}'|} \right) 4\pi G \rho(t', \textbf{x}').
\end{equation}
It is clear that the $4\pi$ terms cancel out and evaluating the time integral over $t'$ (which collapses via the $\delta(t-t')$ function) gives the desired result,
\begin{equation}
    \phi(t, \textbf{x}) = -G \int \frac{\rho(t, \textbf{x}')}{|\textbf{x} - \textbf{x}'|} d^3x'. \label{eq:AppendixPotential}
\end{equation}
This final formula is precisely the standard form of Newton's Law of Universal Gravitation in continuous integral form, successfully verifying the derived Green's function.
 